\documentclass[11pt]{article}
\usepackage[T1]{fontenc}
\usepackage{lmodern}
\usepackage[margin=1in]{geometry}
\usepackage{amsmath,amssymb,amsthm,mathtools}
\usepackage{microtype,booktabs,longtable,array,xcolor,xurl}
\usepackage{enumitem,fancyhdr}
\usepackage[colorlinks=true,linkcolor=blue!40!black,citecolor=blue!40!black,urlcolor=blue!40!black]{hyperref}
\usepackage{bookmark}
\setlength{\headheight}{14pt}
\pagestyle{fancy}\fancyhf{}
\fancyhead[L]{\small Nested harmonic jets and Gamma renormalization}
\fancyhead[R]{\small ProveIt research continuation}
\fancyfoot[C]{\thepage}
\setlength{\parskip}{3pt}
\setlength{\parindent}{1.2em}
\setcounter{tocdepth}{2}
\emergencystretch=2em
\numberwithin{equation}{section}
\newtheorem{theorem}{Theorem}[section]
\newtheorem{proposition}[theorem]{Proposition}
\newtheorem{lemma}[theorem]{Lemma}
\newtheorem{corollary}[theorem]{Corollary}
\theoremstyle{definition}\newtheorem{definition}[theorem]{Definition}
\theoremstyle{remark}\newtheorem{remark}[theorem]{Remark}
\newtheorem{example}[theorem]{Example}
\DeclareMathOperator{\Li}{Li}
\DeclareMathOperator{\FP}{FP}
\DeclareMathOperator{\Log}{Log}
\DeclareMathOperator{\Rea}{Re}
\newcommand{\C}{\mathbb C}\newcommand{\Q}{\mathbb Q}
\newcommand{\N}{\mathbb N}\newcommand{\Z}{\mathbb Z}
\newcommand{\G}{\mathcal G}\newcommand{\F}{\mathfrak F}
\newcommand{\dd}{\,\mathrm d}
\newcommand{\eps}{\varepsilon}
\newcommand{\pin}{\texttt{fc4d3bf80534ad7c901d3b8c9e71baf2df0064ed}}
\newcommand{\src}[1]{\nolinkurl{#1}}
\title{\textbf{Nested Harmonic Jets and Gamma Renormalization}\\[5pt]
\large Exact Stieltjes Identities, Regulator Conversion,\\
Cyclotomic Products, and All-Depth Polylogarithmic Resonance}
\author{Research continuation prepared for Vladimir Reshetnikov\\
\small ProveIt programme \quad\textbullet\quad Prepared with OpenAI ChatGPT}
\date{10 October 2026}
\hypersetup{pdftitle={Nested Harmonic Jets and Gamma Renormalization},pdfauthor={Prepared for Vladimir Reshetnikov with OpenAI ChatGPT},pdfsubject={Exact identities for harmonic sums, polylogarithms and Hurwitz-zeta jets}}
\begin{document}
\maketitle
\begin{abstract}
A single finite harmonic product connects three different operations: taking a
nested-sum generating function, removing a spectral pole, and forming an Abel
transform. We develop their exact relation, keeping their normalizations
separate. A holomorphic two-variable renormalization has the Gamma quotient
$\Gamma(a)/\Gamma(a+z)$ as its spectral value at one. Every spectral jet is
then expressed by a Bell polynomial in generalized Stieltjes constants and
Hurwitz-zeta derivatives. This gives a finite all-depth conversion from
harmonic regularization to diagonal Laurent finite parts, including explicit
correction terms through depth four. Parameter multiplication has a necessary
exponential correction, and a higher-genus cyclotomic identity reduces
fractional-order products to the same Gamma quotient.

For equal-order multiple polylogarithms we classify all depth weights at which
the complete depth-generating function becomes a polynomial in its argument.
At a negative integral depth weight, its first fixed-weight spectral
derivative is elementary. Higher derivatives have explicit logarithmic
harmonic Bell coefficients and reduce to shifted logarithmic Euler sums of
nesting depth at most the derivative order. These results include ordinary
and Euler antiderivative formulas and convergent endpoint identities.
The proofs are analytic and algebraic; 826 exact finite assertions and
1,075 high-precision numerical comparisons accompany them. The latter are
diagnostics, not interval certificates. Classical symmetric-function,
Gamma-product, and regularization results are credited separately; global
priority for the derived families is not asserted.
\end{abstract}
\vspace{4pt}
\noindent\textbf{Scope.} This is an identity-focused continuation, not a claim to
settle the remaining $S_6$ or $S_8$ special-value problems or numerical period
independence. The repository was read at the pinned revision above. No
repository files were modified and no proof-assistant formalization is claimed.
\clearpage
\tableofcontents
\clearpage

\section{Research question, sources, and contribution}
\label{sec:scope}
The integration and differentiation chapters of the ProveIt polylogarithm
manuscript organize single Hurwitz-zeta jets, generalized Stieltjes functions,
and harmonic-number coefficient laws. Its depth chapter and research
programme ask for exact reductions rather than extrapolations from an
integer-relation search. The present continuation asks a compatible but
different question:
\begin{quote}
Can the spectral derivatives of \emph{all nested depths at once} be controlled
by one analytic object, and can special depth weights force exact reductions?
\end{quote}
The answer is affirmative for the equal-order family studied here. There are
two complementary reductions. At the singular endpoint in the spectral
variable, a Gamma-normalized product gives a finite Stieltjes--zeta
coefficient calculus. In the open argument disk, a zero of one finite-product
factor collapses an infinite depth-generating function to a polynomial.
Differentiating near that zero gives the elementary first-order identities
and the higher-order bounded-nesting theorem.

\subsection{What is classical and what is established here}
The finite elementary-symmetric-product identity, Newton's identities,
the Weierstrass product for Gamma, and the relation between symmetric functions
and multiple zeta values are classical. Hoffman develops the latter connection,
including the role of the Gamma genus \cite{Hoffman2008,Hoffman2016}.
Regularization and double-shuffle structures have a substantial independent
theory \cite{IKZ2006,GuoZhang2010}. We do not identify every regularization with
every other one, nor call the reciprocal-Gamma generating function a new
special-function evaluation.

The results proved in this article are the following coordinated statements.
Theorem~\ref{thm:renorm} proves a locally uniform cutoff realization of the
spectral family. Theorems~\ref{thm:jets} and~\ref{thm:regulator} give all-order
jet and regulator-conversion formulas with precise finite coefficient
requirements. Theorems~\ref{thm:multiplication} and~\ref{thm:cyclotomic}
identify the parameter-multiplication correction and the cyclotomic
fractional-order Gamma reduction. Theorems~\ref{thm:truncate},
\ref{thm:firstres}, and~\ref{thm:higherres} classify polynomial truncations,
evaluate every first resonant derivative elementarily, and give every higher
resonant coefficient. Corollary~\ref{cor:depth} specifies the comparison class
in the resulting bounded-nesting reduction. Their derivations are included
in full. An exhaustive literature priority investigation is not part of
this delivery.

\subsection{Pinned repository and targeted reading}
The source revision is
\begin{center}\small\pin.\end{center}
The main source areas are
\src{Analysis/Polylogarithms/docs/manuscript} and \src{docs/incoming}
\cite{ProveIt}. The main manuscript, relevant portions of the depth,
integration, differentiation, and discovery chapters, and the incoming
inventory/intake instructions were inspected. The filenames and README
summaries of all five incoming archives at this snapshot were also examined.
Their complete manuscripts and verification outputs were \emph{not}
exhaustively audited.

Those incoming reports already develop shifted Stieltjes products,
convolution closure, contact terms, ordinary log-Gamma correlations, and
related Lerch limits. The present work does not repackage those bilinear
identities as new results. Its word ``resonance'' refers to the zero
$1-mk^{-s}=0$ in a \emph{depth-generating product}, not to an integer-order
Lerch expansion. The source audit in Section~\ref{sec:audit} states the
remaining limitations and a specific editorial correction.

\section{Finite harmonic products and conventions}
\label{sec:finite}
Fix a real parameter $a>0$, and put $x_n=n+a$. All powers $x_n^s$ use the real
logarithm of the positive base. Spectral differentiation always means
$d/ds$, not differentiation in $a$. Define
\begin{align}
 H_N^{(s)}(a)&=\sum_{n=0}^{N-1}x_n^{-s},\\
 E_{d,N}(s;a)&=\sum_{0\le n_1<\cdots<n_d<N}
                (x_{n_1}\cdots x_{n_d})^{-s},\qquad E_{0,N}=1.
 \label{eq:finiteE}
\end{align}
The function $E_{d,N}$ vanishes for $d>N$. Equal orders make increasing or
decreasing index notation equivalent; the multiple-polylogarithm notation
introduced later will use decreasing indices.

\begin{proposition}[Finite product and Newton recurrence]
For every $N$,
\begin{equation}
 P_N(s,z;a):=\prod_{n=0}^{N-1}(1+z x_n^{-s})
       =\sum_{d=0}^N z^d E_{d,N}(s;a),
 \label{eq:product}
\end{equation}
and
\begin{equation}
 d E_{d,N}(s;a)=\sum_{k=1}^d(-1)^{k-1}
       H_N^{(ks)}(a) E_{d-k,N}(s;a).
 \label{eq:newton}
\end{equation}
\end{proposition}
\begin{proof}
Choose the term $zx_n^{-s}$ from exactly $d$ factors of the product to obtain
\eqref{eq:product}. As a formal power series in $z$,
\[
 z\partial_z\log P_N
 =\sum_{k\ge1}(-1)^{k-1}H_N^{(ks)}(a)z^k.
\]
Multiplication by $P_N$ and comparison of coefficients give
\eqref{eq:newton}. Only finite polynomials are needed for each coefficient.
\end{proof}
Thus generalized harmonic numbers are the power-sum coordinates of the
nested family. For example,
\[
 E_{2,N}=\frac{(H_N^{(s)})^2-H_N^{(2s)}}2,\quad
 E_{3,N}=\frac{(H_N^{(s)})^3-3H_N^{(s)}H_N^{(2s)}+2H_N^{(3s)}}6.
\]
Their spectral derivatives insert logarithms:
\begin{equation}
 \partial_s^r H_N^{(ks)}(a)
   =(-k)^r\sum_{n<N}\frac{\log^r x_n}{x_n^{ks}}.
 \label{eq:harmoniclog}
\end{equation}
The factor $k^r$ in this equation is essential throughout the paper.

We use the Laurent convention
\begin{equation}
 \zeta(1+\eps,a)=\frac1\eps+
     \sum_{r\ge0}\frac{(-1)^r\gamma_r(a)}{r!}\eps^r,
 \qquad \gamma_0(a)=-\psi(a).
 \label{eq:stieltjes}
\end{equation}
The standard Hurwitz identities and their domains are recorded in DLMF
\S25.11 \cite{DLMFzeta}. In particular
$\partial_a\zeta(s,a)=-s\zeta(s+1,a)$ and
$\zeta(k,a)=(-1)^k\psi^{(k-1)}(a)/(k-1)!$ for $k\ge2$.

The complete exponential Bell polynomials $Y_r$ are defined by
\begin{equation}
 \exp\left(\sum_{j\ge1}X_j\frac{t^j}{j!}\right)
       =\sum_{r\ge0}Y_r(X_1,\ldots,X_r)\frac{t^r}{r!}.
 \label{eq:Bell}
\end{equation}
Thus $Y_0=1$, $Y_1=X_1$, $Y_2=X_1^2+X_2$, and
$Y_3=X_1^3+3X_1X_2+X_3$.

\section{A holomorphic cutoff renormalization}
\label{sec:renorm}
Set $X=N+a$ and define the entire-in-$s$ function
\begin{equation}
 C_N(s;a)=\frac{1-X^{1-s}}{s-1},\qquad
 C_N(1;a)=\log X.
 \label{eq:counterterm}
\end{equation}
The endpoint $N+a$, and not a silently substituted endpoint, is part of the
normalization. Its derivatives at one are
\begin{equation}
 \partial_s^r C_N(1;a)=\frac{(-1)^r}{r+1}\log^{r+1}X.
 \label{eq:counterderivative}
\end{equation}
Define
\begin{equation}
 \G_N(s,z;a)=e^{-zC_N(s;a)}P_N(s,z;a).
 \label{eq:GN}
\end{equation}

\begin{theorem}[Cutoff limit and the Gamma slice]
\label{thm:renorm}
For fixed $a>0$, the functions $\G_N$ converge locally uniformly on
$\{\Rea s>1/2\}\times\C$ to a holomorphic function $\G$, entire in $z$,
given by
\begin{equation}
 \G(s,z;a)=
 e^{z(\zeta(s,a)-1/(s-1))}
 \prod_{n\ge0}(1+zx_n^{-s})e^{-zx_n^{-s}}.
 \label{eq:G}
\end{equation}
The apparent singularity at $s=1$ is removable. All fixed mixed $s,z$
derivatives converge locally uniformly. Moreover
\begin{equation}
 \boxed{\G(1,z;a)=\frac{\Gamma(a)}{\Gamma(a+z)}}.
 \label{eq:gammaslice}
\end{equation}
The zeros of $\G(s,\cdot\,;a)$ are exactly $z=-x_n^s$, each simple.
\end{theorem}
\begin{proof}
Euler summation gives, locally uniformly for $\Rea s>0$,
\[
 \zeta(s,X)-\frac{X^{1-s}}{s-1}=O(X^{-\Rea s}),
\]
where the difference is interpreted at one by removability. One direct
justification is to subtract the integral of $(X+t)^{-s}$ from its sum over
nonnegative integers, interval by interval. The resulting series of
differences is bounded on a compact $s$-set by a constant times
$\int_0^\infty(X+t)^{-\Rea s-1}\dd t$, plus its endpoint term. Since
$H_N^{(s)}(a)=\zeta(s,a)-\zeta(s,X)$, this proves
\[
 H_N^{(s)}(a)-C_N(s;a)
 \longrightarrow\zeta(s,a)-\frac1{s-1}.
\]
For a compact set in the stated domain, choose $\sigma>1/2$ below every
real part and a uniform bound for $|z|$. The tail factors satisfy
\[
 \log(1+zx_n^{-s})-zx_n^{-s}=O(x_n^{-2\sigma})
\]
for all sufficiently large $n$, with the logarithm chosen near one.
The summable majorant proves normal convergence of the centered product.
Combining its convergence with that of the linear exponential proves
\eqref{eq:G}. The holomorphic convergence theorem and Cauchy's integral
formula give every fixed derivative assertion.

At one the first exponential is $e^{-z\psi(a)}$. The Weierstrass Gamma
product \cite{DLMFgamma} gives
\[
 \frac{\Gamma(a)}{\Gamma(a+z)}
 =e^{-z\psi(a)}\prod_{n\ge0}
       \left(1+\frac z{x_n}\right)e^{-z/x_n},
\]
proving \eqref{eq:gammaslice}. The exponential never vanishes. Away from
the displayed linear-factor zeros, the normally convergent tail has a
nonzero logarithm and hence is nonzero. Distinct $x_n$ give distinct moduli
$x_n^{\Rea s}$, so the factors have no coincident zeros.
\end{proof}

\begin{corollary}[Harmonic regularized depth coefficients]
\label{cor:harmonicreg}
Write $g_d(a)=[z^d]\Gamma(a)/\Gamma(a+z)$. Then
\begin{equation}
 g_d(a)=\lim_{N\to\infty}
       \sum_{j=0}^d\frac{(-\log(N+a))^j}{j!}E_{d-j,N}(1;a).
 \label{eq:harmonicreg}
\end{equation}
More generally, for every fixed $d,r$, coefficient extraction and
$r$ spectral derivatives commute with the limit in \eqref{eq:GN}.
\end{corollary}
\begin{proof}
Expand the finite exponential in $z$, then apply the local uniform
convergence on a circle about $z=0$. Apply Cauchy's formula also in $s$
for the derivative statement.
\end{proof}
This is a coefficientwise regularization, not an ordinary value assigned to
a divergent nested sum without a subtraction prescription.

\section{The Stieltjes--zeta spectral jet calculus}
\label{sec:jets}
For $z$ initially near zero define the logarithm by $\log\G(s,0;a)=0$, and
put
\begin{equation}
 \Lambda_r(z;a)=\left.\partial_s^r\log\G(s,z;a)\right|_{s=1},
 \qquad F_r(z;a)=\left.\partial_s^r\G(s,z;a)\right|_{s=1}.
 \label{eq:Lambda}
\end{equation}

\begin{theorem}[All spectral jets]
\label{thm:jets}
For $r\ge0$ and $|z|<a$,
\begin{equation}
 \boxed{\Lambda_r(z;a)=(-1)^r z\gamma_r(a)+
       \sum_{k\ge2}(-1)^{k-1}k^{r-1}z^k\zeta^{(r)}(k,a)}.
 \label{eq:logjet}
\end{equation}
Here $\zeta^{(r)}$ differentiates its first variable and $k^{r-1}=1/k$
when $r=0$. For $r\ge1$,
\begin{equation}
 F_r(z;a)=\frac{\Gamma(a)}{\Gamma(a+z)}
          Y_r(\Lambda_1(z;a),\ldots,\Lambda_r(z;a)).
 \label{eq:Belljet}
\end{equation}
Each $F_r$ is entire. The apparent poles on the right of \eqref{eq:Belljet}
cancel. The functions $\Lambda_r$, $r\ge1$, continue meromorphically to the
$z$-plane, with possible poles only at $-a,-a-1,\ldots$.
\end{theorem}
\begin{proof}
For small $z$ the logarithm of \eqref{eq:G} is
\begin{equation}
 \log\G(s,z;a)=z\left(\zeta(s,a)-\frac1{s-1}\right)
    +\sum_{k\ge2}\frac{(-1)^{k-1}z^k}{k}\zeta(ks,a).
 \label{eq:logG}
\end{equation}
Normal convergence in a neighborhood of $s=1$ allows spectral
differentiation. Equation~\eqref{eq:stieltjes} gives the first term, while
$d^r\zeta(ks,a)/ds^r=k^r\zeta^{(r)}(ks,a)$ gives the factor in
\eqref{eq:logjet}. Expanding the exponential in the increment $s-1$
proves \eqref{eq:Belljet}. Entireness of $F_r$ follows from
Theorem~\ref{thm:renorm}, independently of any cancellation manipulation.

For continuation of $\Lambda_r$, separate any finite number of factors from
the product. A spectral derivative of their logarithms is rational in $z$
for $r\ge1$. The remaining differentiated centered logarithmic series
converges normally away from its factor zeros: a fixed compact set has a
summable tail of size $O(x_n^{-2}\log^r x_n)$. This proves the meromorphic
statement and also the cancellation claim in \eqref{eq:Belljet}.
\end{proof}

At order zero, $\Lambda_0=\log\Gamma(a)-\log\Gamma(a+z)$ denotes the
analytic logarithm continued from zero, not the principal logarithm of a
possibly winding Gamma value. The first depth coefficients are
\begin{align}
 g_1&=c,\qquad g_2=\frac{c^2-Z_2}{2},\\
 g_3&=\frac{c^3-3cZ_2+2Z_3}{6},\\
 g_4&=\frac{c^4-6c^2Z_2+3Z_2^2+8cZ_3-6Z_4}{24},
 \label{eq:gtable}
\end{align}
where $c=\gamma_0(a)$ and $Z_k=\zeta(k,a)$. At zero spectral order, the
$Z_k$ may equally be replaced by ordinary polygamma values.

\subsection{An explicit logarithmic-harmonic limit}
Differentiating \eqref{eq:GN} logarithmically before taking the limit gives
\begin{equation}
 \Lambda_1(z;a)=\lim_{N\to\infty}
 \left\{\frac z2\log^2(N+a)
          -\sum_{n<N}\frac{z\log x_n}{x_n+z}\right\}.
 \label{eq:L1limit}
\end{equation}
This identity holds away from the poles and locally uniformly there.
It expresses a Stieltjes--zeta jet as a renormalized generalized harmonic
sum with a shifted denominator. Higher limits follow from
\eqref{eq:counterderivative}; they must differentiate the counterterm too.

\subsection{Exact cancellation at a moving Gamma zero}
\begin{proposition}[Zero-flow identities]
\label{prop:zeroflow}
Let $x=a+m$ with $m\ge0$. Then
\begin{equation}
 F_1(-x;a)=x\log x\,\Gamma(a)(-1)^m m!.
 \label{eq:Froot}
\end{equation}
More generally, with $V_x=\partial_s+(\log x)z\partial_z$,
\begin{equation}
 (V_x^r\G)(s,-x^s;a)=0\qquad(r\ge0,\ \Rea s>1/2).
 \label{eq:zeroflow}
\end{equation}
\end{proposition}
\begin{proof}
The path $z(s)=-x^s$ is a zero of the corresponding linear factor of
\eqref{eq:G}. Its total derivative is precisely $V_x$, proving
\eqref{eq:zeroflow}. At one the $r=1$ case says
$F_1(-x)=x\log x\,\partial_z\G(1,-x)$.
The reciprocal Gamma function has derivative $(-1)^m m!$ at $-m$;
\eqref{eq:gammaslice} gives \eqref{eq:Froot}.
\end{proof}
Thus a zero of the Gamma slice need not be a zero of its spectral derivative.
For instance $x=1$ gives zero because $\log x=0$, whereas a generic $x$
gives the nonzero value in \eqref{eq:Froot}.

\section{Diagonal Laurent finite parts are a different regulator}
\label{sec:regulator}
For $\Rea s>1$ define the diagonal multiple Hurwitz zeta value
\begin{equation}
 Z_d(s;a)=\sum_{0\le n_1<\cdots<n_d}
        (x_{n_1}\cdots x_{n_d})^{-s},\qquad Z_0=1.
 \label{eq:diagZ}
\end{equation}
Newton's identities express every fixed $Z_d$ as a polynomial in
$\zeta(s,a),\ldots,\zeta(ds,a)$, so it is meromorphic in $s$.
Its Laurent finite part at one means the coefficient of $(s-1)^0$ of
\emph{this specific one-variable restriction}.

\begin{theorem}[Exact conversion of the two regulators]
\label{thm:regulator}
As a formal depth series with meromorphic coefficients near $\eps=0$,
\begin{equation}
 \sum_{d\ge0}z^d Z_d(1+\eps;a)
       =e^{z/\eps}\G(1+\eps,z;a).
 \label{eq:regmaster}
\end{equation}
Consequently, for $d\ge0$ and any integer $v$,
\begin{equation}
 [\eps^v]Z_d(1+\eps;a)
   =\sum_{j=0}^d\frac{[z^{d-j}\eps^{v+j}]\G(1+\eps,z;a)}{j!},
 \label{eq:reggeneral}
\end{equation}
where a negative Taylor index on $\G$ means zero. In particular,
\begin{equation}
 \boxed{\FP_{s=1}Z_d(s;a)=
       \sum_{j=0}^d\frac{[z^{d-j}]F_j(z;a)}{(j!)^2}}.
 \label{eq:regFP}
\end{equation}
The finite part belongs to the finite polynomial algebra over $\Q$
generated by
\begin{equation}
 \gamma_0(a),\ldots,\gamma_{d-1}(a),\qquad
 \zeta^{(j)}(k,a)\quad(2\le k\le d,\ 0\le j\le d-k)
 \label{eq:FPatoms}
\end{equation}
for $d\ge1$.
\end{theorem}
\begin{proof}
In the absolutely convergent region the depth-generating function of
\eqref{eq:diagZ} is the uncentered infinite product. Equation~\eqref{eq:G}
shows that it equals $e^{z/(s-1)}\G(s,z;a)$. Alternatively, logarithmic
expansion gives the same identity immediately from Newton's recurrence.
For each fixed depth it is an equality of meromorphic functions, so it
continues to a punctured neighborhood of one.

In a fixed coefficient $z^d$, the exponential contributes only its terms
$j=0,\ldots,d$. Extracting the Laurent coefficient proves
\eqref{eq:reggeneral}. Since $[\eps^j]\G=F_j/j!$, this gives
\eqref{eq:regFP}, including the second factorial. The atom restriction
follows by substituting \eqref{eq:logjet}: a $z^k$ zeta term needs $k$
units of depth, leaving at most $d-k$ spectral derivatives from the pole
factor. A linear Stieltjes term leaves at most $d-1$ derivatives. Products
of such terms add no new generators.
\end{proof}

The first nontrivial formulas, with $c=\gamma_0(a)$, are
\begin{align}
 \FP Z_2&=g_2-\gamma_1(a),\label{eq:FP2}\\
 \FP Z_3&=g_3-c\gamma_1(a)-\zeta'(2,a)+\frac{\gamma_2(a)}4,
 \label{eq:FP3}\\
 \FP Z_4&=g_4+\frac{\zeta(2,a)\gamma_1(a)}2
       -c\zeta'(2,a)-\frac{\zeta''(2,a)}2+\zeta'(3,a)\notag\\
 &\hspace{8mm}-\frac{c^2\gamma_1(a)}2
       +\frac{c\gamma_2(a)}4+\frac{\gamma_1(a)^2}4
       -\frac{\gamma_3(a)}{36}.
 \label{eq:FP4}
\end{align}
For example, the elementary diagonal identity
$Z_2(s;a)=(\zeta(s,a)^2-\zeta(2s,a))/2$ already proves
\eqref{eq:FP2}. Its extra $-\gamma_1(a)$ comes from multiplying the
spectral pole by the next Laurent coefficient. Setting $s=1$ after removing
only the depth-one pole does not reproduce this finite part.

\begin{remark}[Scope of ``regularized'']
Harmonic cutoff regularization \eqref{eq:harmonicreg}, the diagonal Laurent
finite part \eqref{eq:regFP}, and shuffle regularization are not interchangeable
names. The identities above specify the first two directly. They do not
assert that raw diagonal finite-part extraction is a shuffle character.
The general double-shuffle regularization theory \cite{IKZ2006} addresses
additional algebraic compatibility questions. A different finite subtraction
constant changes the harmonic normalization by an exponential in $z$.
\end{remark}

\section{Parameter multiplication and its spectral correction}
\label{sec:multiplication}
\begin{theorem}[Shift and multiplication identities]
\label{thm:multiplication}
For $a>0$, $\Rea s>1/2$, and all $z$,
\begin{equation}
 \G(s,z;a)=(1+z/a^s)\G(s,z;a+1).
 \label{eq:shiftG}
\end{equation}
For every integer $q\ge2$,
\begin{equation}
 \boxed{\prod_{j=0}^{q-1}\G\left(s,z;\frac{a+j}{q}\right)
 =\exp\left(z\frac{q^s-q}{s-1}\right)\G(s,q^sz;a)}.
 \label{eq:multG}
\end{equation}
The quotient in the exponential is removable at one, where its value is
$q\log q$.
\end{theorem}
\begin{proof}
The shift removes the $n=0$ term from the uncentered product; the centered
exponential correction removes exactly the same term, giving
\eqref{eq:shiftG}. For multiplication, use
\[
 \sum_{j=0}^{q-1}\zeta\left(t,\frac{a+j}{q}\right)=q^t\zeta(t,a),
\]
Insert this identity into \eqref{eq:logG}. Terms of degree $k\ge2$ agree
with those of $\log\G(s,q^sz;a)$. The difference in degree one is
$z(q^s-q)/(s-1)$. Exponentiate first near $z=0$ and extend by entireness.
The same proof may be carried out in the absolutely convergent region
before analytic continuation.
\end{proof}
At $s=1$, \eqref{eq:multG} is the Gamma multiplication ratio
\begin{equation}
 \prod_{j<q}\frac{\Gamma((a+j)/q)}{\Gamma((a+j)/q+z)}
       =q^{qz}\frac{\Gamma(a)}{\Gamma(a+qz)}.
 \label{eq:gammamult}
\end{equation}
The classical Gamma multiplication formula is recorded in DLMF \S5.5
\cite{DLMFmult}. Omitting the exponential in \eqref{eq:multG} would already
contradict this specialization.

\begin{corollary}[All-order distribution of logarithmic jets]
\label{cor:jetdist}
Let $\ell=\log q$ and $D=w\partial_w$. Then, as germs at $z=0$,
\begin{align}
 \sum_{j=0}^{q-1}\Lambda_r\left(z;\frac{a+j}{q}\right)
 ={}&\frac{qz\ell^{r+1}}{r+1}\notag\\
 &+\left.\sum_{h=0}^r\binom rh\ell^{r-h}
              D^{r-h}\Lambda_h(w;a)\right|_{w=qz}.
 \label{eq:jetdist}
\end{align}
For $r=1$ this becomes
\begin{equation}
 \sum_{j<q}\Lambda_1\left(z;\frac{a+j}{q}\right)
 =\Lambda_1(qz;a)-qz\ell\,\psi(a+qz)+\frac{qz\ell^2}{2}.
 \label{eq:jetdist1}
\end{equation}
\end{corollary}
\begin{proof}
The $r$-th derivative at one of $(q^s-q)/(s-1)$ is
$q\ell^{r+1}/(r+1)$. Along $w=q^sz$, differentiation of a function of
$(s,w)$ acts as $\partial_s+\ell D$; the two operators commute.
Expanding its $r$-th power gives \eqref{eq:jetdist}. Finally
$D\Lambda_0(w;a)=-w\psi(a+w)$ gives \eqref{eq:jetdist1}.
\end{proof}
This identity links Stieltjes jets, spectral zeta derivatives, and a
polygamma argument shift without any numerical relation search. Further
$z$ derivatives of \eqref{eq:jetdist1} produce explicit higher polygamma
terms and remain valid meromorphically away from the displayed poles.

\section{Higher-genus products and a cyclotomic Gamma identity}
\label{sec:cyclotomic}
The half-plane in Theorem~\ref{thm:renorm} is tied to a quadratic tail. To
cross it, higher power sums must also be subtracted. This can be done
without introducing an unspecified regularization.

\begin{definition}
For an integer $M\ge1$ define
\begin{equation}
 \G_N^{[M]}(s,z;a)=P_N(s,z;a)
  \exp\left(-\sum_{k=1}^M\frac{(-1)^{k-1}z^k}{k}C_N(ks;a)\right).
 \label{eq:genuscutoff}
\end{equation}
Thus $\G_N^{[1]}=\G_N$.
\end{definition}

\begin{theorem}[Higher-genus limit and cyclotomic reduction]
\label{thm:cyclotomic}
The functions in \eqref{eq:genuscutoff} converge locally uniformly for
$\Rea s>1/(M+1)$ and all $z$, to
\begin{align}
 \G^{[M]}(s,z;a)
 ={}&\exp\left(\sum_{k=1}^M\frac{(-1)^{k-1}z^k}{k}
            \left(\zeta(ks,a)-\frac1{ks-1}\right)\right)\notag\\
 &\times\prod_{n\ge0}(1+zx_n^{-s})
       \exp\left(\sum_{k=1}^M\frac{(-1)^kz^k}{kx_n^{ks}}\right).
 \label{eq:higherG}
\end{align}
Every apparent pole at $s=1/k$, $1\le k\le M$, is removed in the
specified combination. Put $\omega=e^{2\pi i/m}$ and
$u=(-1)^{m+1}z^m$. For $m\ge1$ and $\Rea s>1/(m+1)$,
\begin{equation}
 \prod_{j=0}^{m-1}\G^{[m]}(s,\omega^jz;a)=\G(ms,u;a).
 \label{eq:cyclogen}
\end{equation}
In particular,
\begin{equation}
 \boxed{\prod_{j=0}^{m-1}\G^{[m]}(1/m,\omega^jz;a)
       =\frac{\Gamma(a)}{\Gamma(a+(-1)^{m+1}z^m)}}.
 \label{eq:cyclogamma}
\end{equation}
For logarithms continued from $z=0$, every $r\ge0$ also satisfies
\begin{equation}
 \left.\sum_{j=0}^{m-1}\partial_s^r
       \log\G^{[m]}(s,\omega^jz;a)\right|_{s=1/m}
       =m^r\Lambda_r((-1)^{m+1}z^m;a).
 \label{eq:cyclojets}
\end{equation}
\end{theorem}
\begin{proof}
Separate the first $M$ terms of the logarithmic expansion of each factor.
The remaining tail is $O(x_n^{-(M+1)\sigma})$ on a compact set with
$\Rea s\ge\sigma>1/(M+1)$. For each $k\le M$, the Euler-summation
argument in Theorem~\ref{thm:renorm} applies to $ks$, and proves convergence
of $H_N^{(ks)}-C_N(ks)$ to its pole-subtracted zeta value. This proves
\eqref{eq:higherG}, holomorphy, and convergence of all fixed derivatives.

The cyclotomic assertion is already an exact finite-cutoff identity.
Indeed
\[
 \prod_{j=0}^{m-1}(1+\omega^jt)=1-(-t)^m
             =1+(-1)^{m+1}t^m,
\]
and $\sum_{j<m}\omega^{jk}=0$ unless $m$ divides $k$. In the counterterm
sum with $1\le k\le m$, only $k=m$ remains; it is precisely
$-uC_N(ms;a)$. Hence
\[
 \prod_{j<m}\G_N^{[m]}(s,\omega^jz;a)=\G_N(ms,u;a).
\]
Both sides have locally uniform limits in the stated half-plane, since
$m\Rea s>m/(m+1)\ge1/2$. Taking the limit proves \eqref{eq:cyclogen}.
The Gamma slice gives \eqref{eq:cyclogamma}. Logarithmic differentiation of
\eqref{eq:cyclogen}, on a neighborhood free of zeros, gives
\eqref{eq:cyclojets}; each spectral derivative contributes a factor $m$.
\end{proof}

For example the half-order pair has the exact identity
\begin{equation}
 \G^{[2]}(1/2,z;a)\G^{[2]}(1/2,-z;a)
      =\frac{\Gamma(a)}{\Gamma(a-z^2)}.
 \label{eq:halforder}
\end{equation}
At third order the corresponding three rotated factors give
$\Gamma(a)/\Gamma(a+z^3)$. Their individual jets involve pole-subtracted
Hurwitz values at fractional spectral arguments; the cyclotomic sum removes
all components except the ordinary Stieltjes--zeta jet on the right of
\eqref{eq:cyclojets}. No choice of a branch of $z^{1/m}$ is needed.

\section{Polylogarithmic Euler primitives and ordinary antiderivatives}
\label{sec:primitives}
The spectral coefficients admit a second summation direction. This turns
Euler integration in the depth variable into a shift in an ordinary
polylogarithm's order.

\begin{theorem}[A complete normalized primitive ladder]
\label{thm:primitives}
For integers $r,q\ge0$, initially for $|z|<a$, put
\begin{equation}
 A_{r,q}(z;a)=(-1)^rz\gamma_r(a)+
   \sum_{k\ge2}(-1)^{k-1}k^{r-q-1}z^k\zeta^{(r)}(k,a).
 \label{eq:Aseries}
\end{equation}
Then $A_{r,0}=\Lambda_r$, $A_{r,q}(0;a)=0$, and
\begin{equation}
 z\partial_z A_{r,q+1}(z;a)=A_{r,q}(z;a).
 \label{eq:Eulerladder}
\end{equation}
An exact normally convergent continuation to
$\Omega_a=\C\setminus(-\infty,-a]$ is
\begin{equation}
 \boxed{A_{r,q}(z;a)=(-1)^r\left[
 z\gamma_r(a)-\sum_{n\ge0}\log^r x_n
       \left(\Li_{q+1-r}(-z/x_n)+\frac z{x_n}\right)\right]}.
 \label{eq:Apolylog}
\end{equation}
For $r=0$ the factor $\log^0 x_n$ is defined as one, including $x_n=1$.
The polylogarithms use their principal continuation on their usual slit
plane.
\end{theorem}
\begin{proof}
Apply $z\partial_z$ termwise in \eqref{eq:Aseries}, multiplying a degree
$k$ term by $k$. This proves \eqref{eq:Eulerladder}. For $k\ge2$ the
absolutely convergent Dirichlet representation gives
\[
 \zeta^{(r)}(k,a)=(-1)^r\sum_{n\ge0}\frac{\log^r x_n}{x_n^k}.
\]
Substitute it into \eqref{eq:Aseries} and sum first over $k$ using the
polylogarithm series \cite{DLMFpolylog}. The removed linear term is exactly
$z/x_n$, and the result is \eqref{eq:Apolylog}.

For any compact subset of $\Omega_a$, sufficiently large $x_n$ make the
bracket $O(|z|^2/x_n^2)$, uniformly on that compact set. Its product with
$\log^r x_n$ is summable. The finitely many remaining polylogarithms are
holomorphic there; their individual cuts start at $-x_n$ and are contained
in the excluded ray. This proves normal convergence and analytic
continuation. Equation~\eqref{eq:Eulerladder} consequently continues too.
\end{proof}
This is a normalized ladder, rather than an ambiguous notation for negative
polygamma. Every Euler primitive vanishes at zero. Negative integer orders
of the polylogarithm in \eqref{eq:Apolylog} are rational functions.
In particular,
\begin{align}
 \Lambda_1(z;a)&=-z\gamma_1(a)
   +z^2\sum_{n\ge0}\frac{\log x_n}{x_n(x_n+z)},
 \label{eq:L1rational}\\
 \Lambda_2(z;a)&=z\gamma_2(a)
   -z^2\sum_{n\ge0}\frac{(\log x_n)^2(2x_n+z)}{x_n(x_n+z)^2}.
 \label{eq:L2rational}
\end{align}
These formulas also give the meromorphic continuation asserted in
Theorem~\ref{thm:jets} for the first two orders.

\begin{corollary}[Two ordinary antiderivatives]
\label{cor:ordinary}
For $z\in\Omega_a$, with path integration from zero within that domain,
\begin{align}
 \int_0^z\Lambda_1(u;a)\dd u
 ={}&-\frac{\gamma_1(a)z^2}{2}\notag\\
 &+\sum_{n\ge0}\log x_n\left[
 \frac{z^2}{2x_n}-z+x_n\Log\left(1+\frac z{x_n}\right)\right],
 \label{eq:ordinaryL1}\\
 \int_0^z\Lambda_0(u;a)\dd u
 ={}&z\log\Gamma(a)-\zeta'(-1,a+z)+\zeta'(-1,a)\notag\\
 &+az+\frac{z^2-z}{2}-\frac z2\log(2\pi).
 \label{eq:ordinaryL0}
\end{align}
The series in \eqref{eq:ordinaryL1} is normally convergent.
\end{corollary}
\begin{proof}
The bracket in \eqref{eq:ordinaryL1} is the integral of
$u^2/(x_n(x_n+u))$, and its tail starts with $z^3/(3x_n^2)$.
This justifies integration term by term in \eqref{eq:L1rational}.
For \eqref{eq:ordinaryL0}, differentiate
$\partial_a\zeta(s,a)=-s\zeta(s+1,a)$ in $s$ at $s=-1$ and use
$\zeta'(0,a)=\log\Gamma(a)-\tfrac12\log(2\pi)$ and
$\zeta(0,a)=\tfrac12-a$ \cite{DLMFzeta}. The derivative of the asserted
right side is $\Lambda_0(z;a)$ and its value at zero is zero.
\end{proof}
Formula~\eqref{eq:ordinaryL0} is the classical log-Gamma primitive in the
present normalization, included as a compatibility check rather than a
new evaluation. Formula~\eqref{eq:ordinaryL1} is its spectral-jet companion.

\section{The complete depth-generating polylogarithm}
\label{sec:polygen}
For $|w|<1$ use decreasing indices and define
\begin{equation}
 \Li_{\{s\}^d}(w,1,\ldots,1)
 =\sum_{n_1>\cdots>n_d\ge1}
       \frac{w^{n_1}}{(n_1\cdots n_d)^s}.
 \label{eq:multipleLi}
\end{equation}
In this notation $\partial_s$ differentiates \emph{all $d$ orders
simultaneously along the diagonal}. It is not a derivative of the first
order with the other orders fixed.

\begin{theorem}[Depth product, recurrence, and Abel transform]
\label{thm:polygen}
For $\Rea s>0$, $z\in\C$, and $|w|<1$,
\begin{align}
 \F(s,z,w)&=1+\sum_{d\ge1}z^d\Li_{\{s\}^d}(w,1,\ldots,1)
            =1+\sum_{n\ge1}b_n(s,z)w^n,\label{eq:Fdef}\\
 b_n(s,z)&=\frac z{n^s}\prod_{k=1}^{n-1}\left(1+\frac z{k^s}\right).
 \label{eq:bn}
\end{align}
The depth and coefficient series converge absolutely and locally uniformly
in this domain, as do all fixed parameter derivatives. With $b_0=1$,
\begin{equation}
 n^s b_n=((n-1)^s+z)b_{n-1},\qquad n\ge1,
 \label{eq:brec}
\end{equation}
where $0^s=0$. There is also the exact Abel-transform identity
\begin{equation}
 \F(s,z,w)=(1-w)\sum_{N\ge0}P_N(s,z;1)w^N.
 \label{eq:Abel}
\end{equation}
On the slice $s=1$,
\begin{equation}
 \F(1,z,w)=(1-w)^{-z}.
 \label{eq:Fones}
\end{equation}
\end{theorem}
\begin{proof}
Fix the largest index $n_1=n$. The remaining factors enumerate elementary
symmetric functions of $1^{-s},\ldots,(n-1)^{-s}$. Summing their depths
proves \eqref{eq:bn}. On a compact set let $\sigma>0$ bound the real parts
from below, $|z|\le R$, and $|w|\le\rho<1$. The absolute contribution
at largest index $n$ is at most
\[
 \rho^n Rn^{-\sigma}\exp\left(R\sum_{k<n}k^{-\sigma}\right).
\]
The exponent is $o(n)$ for every $\sigma>0$, so this is summable.
It bounds the depth series as well as the coefficient series and proves
joint local uniform convergence. Holomorphy and Cauchy's formula then
justify all fixed derivatives and regroupings.

Cancel the common product factors to obtain \eqref{eq:brec}, or verify it
without division by multiplying the products. Since
$b_n=P_n(s,z;1)-P_{n-1}(s,z;1)$, summation by parts proves
\eqref{eq:Abel}. At $s=1$, the recurrence is that of the binomial expansion:
$b_n=(z)_n/n!$, proving \eqref{eq:Fones}.
\end{proof}

For $\Rea s>1$ the endpoint series converges absolutely and gives
\begin{equation}
 \G(s,z;1)=e^{-z/(s-1)}\F(s,z,1).
 \label{eq:FtoG}
\end{equation}
Thus the two principal objects of this article arise from the same finite
products; they are not unrelated special-function constructions.
Nevertheless their normalizations at $s=1$ differ:
\begin{equation}
 \lim_{N\to\infty}(N+1)^{-z}P_N(1,z;1)=\frac1{\Gamma(1+z)},
 \qquad (1-w)^z\F(1,z,w)=1.
 \label{eq:AbelGamma}
\end{equation}
The Gamma factor is the exact distinction between this harmonic cutoff and
this Abel normalization. One may not replace one by the other without it.

\begin{remark}[Why the open spectral half-plane matters]
At $s=0$, direct geometric summation, in its initial convergence region,
gives
\begin{equation}
 \F(0,z,w)=\frac{1-w}{1-(1+z)w}.
 \label{eq:s0boundary}
\end{equation}
For fixed nonzero $w$ this has a pole in $z$. Thus the entireness-in-depth
claim of Theorem~\ref{thm:polygen} cannot simply be extended to $s=0$.
\end{remark}

For later use define $\theta=w\partial_w$, and define $\theta^s$ on a power
series by multiplying its positive-index coefficient by $n^s$ and killing
its constant coefficient. Equation~\eqref{eq:brec} is equivalently
\begin{equation}
 (1-w)\theta^s\F(s,z,w)=zw\F(s,z,w).
 \label{eq:theta}
\end{equation}
This is a coefficient-defined operator identity. No fractional-power
branch for $w$ is being assumed. Its spectral derivative uses the similarly
diagonal multiplier $\log n$ on positive indices.

\section{Polynomial truncations and elementary first spectral derivatives}
\label{sec:resonance}
\begin{theorem}[Complete classification of polynomial truncations]
\label{thm:truncate}
Fix $\Rea s>0$. The function $\F(s,z,w)$ is a polynomial in $w$ if and only
if $z=0$ or $z=-m^s$ for some positive integer $m$. In the latter case its
degree is exactly $m$, and
\begin{equation}
 P_m^*(s,w):=\F(s,-m^s,w)
 =1-\sum_{n=1}^m\left(\frac mn\right)^s
    \prod_{k=1}^{n-1}\left(1-\left(\frac mk\right)^s\right)w^n.
 \label{eq:respoly}
\end{equation}
For every integer $r\ge1$ and $|w|<1$,
\begin{equation}
 \boxed{\sum_{d\ge1}(-m)^d
   \left.(\partial_s+d\log m)^r
       \Li_{\{s\}^d}(w,1,\ldots,1)\right|_{s=1}
       =\left.\partial_s^rP_m^*(s,w)\right|_{s=1}}.
 \label{eq:movingres}
\end{equation}
The right side is a finite polynomial in $w$, with coefficients in the
algebra over $\Q$ generated by $\log 2,\ldots,\log m$.
\end{theorem}
\begin{proof}
At $z=-m^s$, the factor indexed by $k=m$ in \eqref{eq:bn} kills every
$b_n$ with $n>m$. The coefficient $b_m$ is nonzero, since
$k^s\ne m^s$ for $k<m$; their moduli are distinct. This proves
\eqref{eq:respoly} and its degree. Conversely, if a nonconstant sequence
terminates, some $b_n$ with $n\ge2$ vanishes. Its product expression and
$z\ne0$ force $z=-k^s$ for one of the factors. The same modulus argument
makes this $k$ unique.

Along $z(s)=-m^s$, the derivative of the depth coefficient $z(s)^d$ acts
as multiplication by $d\log m$. Local uniform convergence from
Theorem~\ref{thm:polygen} allows differentiation along this curve. This
proves \eqref{eq:movingres}. At $s=1$, derivatives of the finitely many
positive-base powers in \eqref{eq:respoly} have the asserted logarithmic
coefficients.
\end{proof}

The moving-depth derivative in \eqref{eq:movingres} is not the fixed-depth
spectral derivative. The distinction can be resolved exactly at first order.
Put
\begin{equation}
 T_m(w)=\left.\partial_s\F(s,-m,w)\right|_{s=1},\qquad
 Q_m(w)=\left.\partial_sP_m^*(s,w)\right|_{s=1}.
 \label{eq:TQ}
\end{equation}

\begin{theorem}[Elementary first resonant derivative]
\label{thm:firstres}
For every $m\ge1$ and $|w|<1$,
\begin{equation}
 \boxed{T_m(w)=Q_m(w)-m\log m\,(1-w)^m\Log(1-w)}.
 \label{eq:Tclosed}
\end{equation}
Here $Q_m$ is an explicitly computable polynomial of degree at most $m$.
Equivalently,
\begin{equation}
 T_m(w)=-(1-w)^m\sum_{n=1}^m(-1)^n\binom mn n\log n
            \int_0^w\frac{t^{n-1}}{(1-t)^m}\dd t.
 \label{eq:Tintegral}
\end{equation}
All integrals in \eqref{eq:Tintegral} are elementary. In particular,
\begin{equation}
 \boxed{\sum_{d\ge1}(-2)^d
   \left.\partial_s\Li_{\{s\}^d}(w,1,\ldots,1)\right|_{s=1}
 =-2\log2\left[w(1-w)+(1-w)^2\Log(1-w)\right]}.
 \label{eq:T2}
\end{equation}
\end{theorem}
\begin{proof}
Differentiate $P_m^*(s,w)=\F(s,-m^s,w)$ at one. Since
$\partial_z\F(1,z,w)=-\Log(1-w)(1-w)^{-z}$, the chain rule gives
$Q_m=T_m+m\log m(1-w)^m\Log(1-w)$, proving \eqref{eq:Tclosed}.
This already proves elementary reducibility: $Q_m$ is obtained by
differentiating the finite expression \eqref{eq:respoly}.

For an independent constructive expression, differentiate
\eqref{eq:theta} at $s=1,z=-m$, where $\F=(1-w)^m$.
The result, after division by $w(1-w)$ and analytic extension at zero, is
\begin{equation}
 T_m'(w)+\frac m{1-w}T_m(w)
   =-\sum_{n=1}^m(-1)^n\binom mn n\log n\,w^{n-1},
 \qquad T_m(0)=0.
 \label{eq:Tode}
\end{equation}
The integrating factor $(1-w)^{-m}$ gives \eqref{eq:Tintegral}.
Finally
$P_2^*(s,w)=1-2^s w+(2^s-1)w^2$, so
$Q_2=-2\log2\,w(1-w)$; insert this into \eqref{eq:Tclosed}.
\end{proof}

An explicit formula for the elementary integral is
\begin{align}
 \int_0^w\frac{t^{n-1}}{(1-t)^m}\dd t
 ={}&\sum_{\substack{0\le j\le n-1\\j\ne m-1}}
 (-1)^j\binom{n-1}{j}
 \frac{1-(1-w)^{j-m+1}}{j-m+1}\notag\\
 &+\mathbf1_{n=m}(-1)^m\Log(1-w).
 \label{eq:Iexplicit}
\end{align}
Substituting this finite sum into \eqref{eq:Tintegral} gives a direct
rational-coefficient algorithm for $Q_m$. Two further examples are
\begin{align}
 Q_3(w)={}&-3w\log3+
   \left(\frac{15}{2}\log3-3\log2\right)w^2
   +\left(3\log2-\frac92\log3\right)w^3,\label{eq:Q3}\\
 Q_4(w)={}&-8w\log2+22w^2\log2
  +\left(4\log3-\frac{80}{3}\log2\right)w^3\notag\\
 &+\left(\frac{38}{3}\log2-4\log3\right)w^4.
 \label{eq:Q4}
\end{align}
For $m=1$, $\F(s,-1,w)=1-w$ for every $\Rea s>0$, so every positive
spectral derivative is identically zero.

\begin{corollary}[Elementary coefficient certificate]
\label{cor:Tcoeff}
The coefficients of $T_m$ are
\begin{equation}
 [w^n]T_m=
 \begin{cases}
 (-1)^n\binom mn\left[-\log n+
       m\displaystyle\sum_{k=1}^{n-1}\frac{\log k}{k-m}\right],&1\le n\le m,\\[5pt]
 \displaystyle\frac{(-1)^m m\log m}{(n-m)\binom nm},&n>m.
 \end{cases}
 \label{eq:Tcoeff}
\end{equation}
\end{corollary}
\begin{proof}
For $n\le m$ differentiate the nonvanishing product in \eqref{eq:bn}
logarithmically. For $n>m$, only the derivative of its zero factor
$1-m^{1-s}$ contributes. The remaining rational product at one is
$(-1)^m m/((n-m)\binom nm)$, yielding the second line.
\end{proof}
These formulas give a coefficientwise proof check without evaluating a
multiple polylogarithm numerically. They also show why an infinite sum over
depths can simplify even when its separate depth components look complicated.

\section{All higher resonant jets and bounded nesting depth}
\label{sec:higher}
The first derivative is not an isolated coincidence. Its zero-factor
calculation yields a finite Bell formula at every spectral order.
For $m,n\ge1$ and $j\ge1$ define
\begin{equation}
 L_{j,n}^{(m)}=-\mathbf1_{j=1}\log n+
 (-1)^{j+1}\sum_{\substack{1\le k<n\\k\ne m}}
        (\log k)^j\Li_{1-j}(m/k).
 \label{eq:Lmn}
\end{equation}
The polylogarithms here have nonpositive integer orders, so they are
\emph{rational functions} of $m/k$; no continuation across a logarithmic cut
is involved. The excluded value $k=m$ is exactly the isolated zero factor.
For example,
\begin{align}
 L_{1,n}^{(m)}&=-\log n+m\sum_{\substack{k<n\\k\ne m}}
                                   \frac{\log k}{k-m},\\
 L_{2,n}^{(m)}&=-m\sum_{\substack{k<n\\k\ne m}}
                                   \frac{k(\log k)^2}{(k-m)^2}.
 \label{eq:Lfirsttwo}
\end{align}
Write $Y_r(L)$ for the Bell polynomial in
$L_{1,n}^{(m)},\ldots,L_{r,n}^{(m)}$ with these $m,n$ fixed.

\begin{theorem}[Every fixed-weight spectral coefficient]
\label{thm:higherres}
Let
\[
 U_r^{(m)}(w)=\left.\partial_s^r\F(s,-m,w)\right|_{s=1}
            =\mathbf1_{r=0}+\sum_{n\ge1}c_{r,n}^{(m)}w^n.
\]
For every $r\ge0$,
\begin{equation}
 c_{r,n}^{(m)}=
 \begin{cases}
 (-1)^n\binom mn Y_r(L),&1\le n\le m,\\[4pt]
 \displaystyle\frac{(-1)^m m}{(n-m)\binom nm}
   \sum_{j=1}^r(-1)^{j+1}\binom rj(\log m)^jY_{r-j}(L),&n>m.
 \end{cases}
 \label{eq:allcoeff}
\end{equation}
The empty sum on the second line is zero. The coefficient series converges
locally uniformly in $|w|<1$ and represents the indicated analytic
spectral derivative.
\end{theorem}
\begin{proof}
For $n\le m$, none of the factors in $b_n(s,-m)$ vanishes at one. Choose
the analytic logarithm of $b_n(s,-m)/b_n(1,-m)$ near one. The derivatives
of a nonzero factor are
\begin{equation}
 \left.\partial_s^j\log(1-mk^{-s})\right|_{s=1}
       =(-1)^{j+1}(\log k)^j\Li_{1-j}(m/k),\quad k\ne m.
 \label{eq:factorjet}
\end{equation}
For $|mk^{-s}|<1$ this follows from the logarithmic series. For the other
finitely many $k$, it follows by differentiating the rational first
logarithmic derivative; equivalently both sides are the same rational
function of $mk^{-s}$. The term $n^{-s}$ contributes only $-\log n$ to
the first logarithmic derivative. Thus all logarithmic derivatives are
exactly \eqref{eq:Lmn}. Exponentiating their Taylor series with
\eqref{eq:Bell} proves the first line of \eqref{eq:allcoeff}.

For $n>m$, factor the coefficient as
\begin{equation}
 b_n(s,-m)=(1-m^{1-s})R_n(s),\qquad
 R_n(s)=-mn^{-s}\prod_{\substack{k<n\\k\ne m}}(1-mk^{-s}).
 \label{eq:isolatedzero}
\end{equation}
Now $R_n(1)=(-1)^m m/((n-m)\binom nm)$ and its logarithmic derivatives
are again $L_{j,n}^{(m)}$. Also
\[
 \left.\partial_s^j(1-m^{1-s})\right|_{s=1}
        =(-1)^{j+1}(\log m)^j\quad(j\ge1),
\]
while its value is zero. Leibniz's rule and the Bell expansion for $R_n$
give the second line. Local uniform convergence and identification with
$U_r^{(m)}$ follow from Theorem~\ref{thm:polygen}.
\end{proof}

The first three tail polynomials multiplying $R_n(1)$ are
\begin{align}
 B_1&=\ell,\\
 B_2&=2\ell L_1-\ell^2,\\
 B_3&=3\ell(L_1^2+L_2)-3\ell^2L_1+\ell^3,
 \label{eq:tailBell}
\end{align}
where $\ell=\log m$. The vanishing factor lowers the number of harmonic
sums that can occur by one. This is the mechanism behind the depth bound.

\begin{definition}[Comparison class for the depth reduction]
A shifted logarithmic Euler sum of nesting depth $d+1$ is a convergent sum
of the form
\begin{equation}
 \sum_{n>k_1>\cdots>k_d\ge1}
 w^n R_0(n)(\log n)^{p_0}
       \prod_{i=1}^d R_i(k_i)(\log k_i)^{p_i},
 \label{eq:comparisonclass}
\end{equation}
with finitely many stated index exclusions, nonnegative integers $p_i$,
and rational functions $R_i$ with rational coefficients. In the present
application the exclusions are $n>m$ and $k_i\ne m$, and the poles are
integer shifts arising from $0,1,\ldots,m$. Finite polynomials are included
as depth-zero terms.
\end{definition}

\begin{corollary}[All-depth to bounded-nesting reduction]
\label{cor:depth}
For fixed $m\ge1$ and $r\ge1$, the function
\begin{equation}
 \sum_{d\ge1}(-m)^d
       \left.\partial_s^r\Li_{\{s\}^d}(w,1,\ldots,1)\right|_{s=1}
 \label{eq:depthcollapse}
\end{equation}
is a finite linear combination of a polynomial and shifted logarithmic
Euler sums of nesting depth at most $r$, with coefficients in
$\Q[\log m]$. The coefficients and summands are explicitly generated by
\eqref{eq:allcoeff}. This is a representation theorem in the comparison
class \eqref{eq:comparisonclass}, not a minimal-depth or arithmetic
independence theorem for ordinary multiple polylogarithms.
\end{corollary}
\begin{proof}
Every $Y_{r-j}$ on the second line of \eqref{eq:allcoeff} is a polynomial
with at most $r-j\le r-1$ factors of the single harmonic sums in
\eqref{eq:Lmn}; factors $\log n$ do not introduce another summation index.
Expand their products by partitioning the index tuples into strict order
patterns and equality blocks. An equality block multiplies rational
kernels and adds logarithmic exponents, staying in the stated class.
This is the finite stuffle decomposition and produces at most $r-1$
inner nested indices. The outer index $n$ raises the total nesting depth
to at most $r$. Rationality of every kernel follows from the nonpositive
polylogarithm orders in \eqref{eq:Lmn}. Absolute convergence for $|w|<1$
justifies each finite rearrangement.
\end{proof}

\begin{example}[The second derivative at depth weight $-2$]
Let $\ell=\log2$ and
\[
 M_n=2\sum_{k=3}^{n-1}\frac{\log k}{k-2}-\log n.
\]
The theorem gives the explicit depth-two representation
\begin{equation}
 U_2^{(2)}(w)=\ell^2w^2+
   4\sum_{n=3}^{\infty}
       \frac{w^n(2\ell M_n-\ell^2)}{n(n-1)(n-2)}.
 \label{eq:U22}
\end{equation}
This is an exact identity, not a claim that the second derivative is
elementary. The first derivative's elementary form does not imply the
same conclusion at higher orders.
\end{example}

\begin{corollary}[Convergent endpoint cancellation]
\label{cor:endpoint}
For every $m\ge1$ and $r\ge0$, the coefficient series for $U_r^{(m)}$
converges absolutely at $w=1$, and
\begin{equation}
 U_r^{(m)}(1)=0.
 \label{eq:endpoint}
\end{equation}
Here the endpoint is defined by the \emph{coefficient series in $n$}, or
equivalently its Abel limit. It is not a termwise substitution of $w=1$
into the generally divergent individual-depth multiple polylogarithms.
\end{corollary}
\begin{proof}
For large $k$, $\Li_{1-j}(m/k)=O(1/k)$, so
$L_{j,n}^{(m)}=O(1+\log^{j+1}n)$ at fixed $m,j$. Bell monomials of
weight $r-j$ grow at most as a fixed multiple of
$1+\log^{2(r-j)}n$. Since $R_n(1)=O(n^{-m-1})$, the tail in
\eqref{eq:allcoeff} is absolutely summable; the case $r=0$ is a polynomial.

For every finite $N$, the coefficient sum through $n=N$ is the $r$-th
spectral derivative at one of $P_N(s,-m;1)$ (including its constant for
$r=0$). Write this product as
\[
 P_N(s,-m;1)=e^{-mC_N(s;1)}\G_N(s,-m;1).
\]
The derivatives of $\G_N$ of each fixed order at one are bounded uniformly
in $N$ by Theorem~\ref{thm:renorm}. The exponential contributes
$(N+1)^{-m}$ times a polynomial in $\log(N+1)$ of bounded degree when
differentiated a fixed number of times, by \eqref{eq:counterderivative}.
Their product therefore tends to zero. Absolute convergence identifies
this limit with \eqref{eq:endpoint}.
\end{proof}
For instance \eqref{eq:U22} yields the convergent identity
\begin{equation}
 \sum_{n=3}^{\infty}
 \frac{2\sum_{k=3}^{n-1}\log k/(k-2)-\log n}{n(n-1)(n-2)}=0.
 \label{eq:exampleendpoint}
\end{equation}
This low-order example also follows by interchanging its two absolutely
convergent sums and telescoping the rational tail. The all-order theorem
explains the entire Bell family of such cancellations.

\section{Verification, source audit, and integration scope}
\label{sec:audit}
\subsection{Reproducible evidence and its limits}
The analytic theorems are justified by the proofs above. The delivered
scripts are implementation checks and transparent finite certificates.
They do not turn an approximate transcendental equality into a proof.
The actual executed suites are summarized in Table~\ref{tab:verification}.

\begin{table}[htbp]
\centering\small
\begin{tabular}{@{}p{.30\textwidth}p{.12\textwidth}p{.49\textwidth}@{}}
\toprule
Suite & Checks & Executed scope\\
\midrule
Exact algebra & 826 & Newton products; diagonal Laurent conversion through depth four and Laurent index two; first resonances $1\le m\le12$, $1\le n\le48$; rational negative-order polylogs through order eight.\\[3pt]
Whole-function numerics & 99 & Gamma product, shifts, multiplication, Stieltjes--zeta jets, polylog primitives, zero values, first resonances, moving truncations, and Cauchy finite-part extraction.\\[3pt]
Higher-jet coefficients & 960 & Two independent coefficient constructions for $1\le m\le6$, $1\le n\le32$, and spectral orders zero through four.\\[3pt]
Cyclotomic products & 16 & Genera two through five, two positive Hurwitz parameters, and resonant and nearby complex spectral points.\\
\bottomrule
\end{tabular}
\caption{Executed verification. The last three suites use 60 decimal-digit
floating-point arithmetic; they are not interval certificates.}
\label{tab:verification}
\end{table}

The exact checker uses SymPy 1.14.0 and exact rational arithmetic. In the
first-resonance coefficient comparison, the symbols for $\log2,\log3,\ldots$
are kept formally independent; the check does not rely on numerical
approximations or unproved independence of the actual logarithms.
The finite Laurent arrays are generated by Newton's recurrence, separately
with and without the spectral pole, before their coefficient conversion is
compared. This tests in particular both factorials in \eqref{eq:regFP}.

The numerical suites use Python 3.13.5 and mpmath 1.3.0. The largest scaled
residual among the 99 whole-function comparisons was approximately
$3.35\times10^{-50}$; the largest absolute residual was approximately
$1.13\times10^{-48}$. The higher-jet and cyclotomic suites had largest scaled
residuals approximately $7.16\times10^{-60}$ and $1.02\times10^{-60}$,
respectively. The JSON outputs retain the individual parameters, acceptance
thresholds, and observed residuals. These figures describe the executed
arithmetic, not a certified error interval.

Independent representations are used where practical. Gamma ratios are
compared with finite centered products and Hurwitz-zeta tails. Spectral
zeta series are compared with finite polylogarithmic heads and separately
accelerated tails. First resonant derivatives are obtained by propagating
the differentiated coefficient recurrence and compared with a derivative
of the finite moving-root polynomial. Higher coefficients compare that
recurrence with the zero-factor Bell formula. Diagonal finite parts are
also extracted by a finite-radius Cauchy average of the raw Newton
polynomial, rather than obtained only by subtracting two truncated formal
expressions.

For a reproducible numerical product evaluation one may split at an integer
$N$ and use, for example,
\begin{align}
 \G(s,z;a)={}&e^{z(\zeta(s,a)-1/(s-1))}
 \prod_{n<N}(1+zx_n^{-s})e^{-zx_n^{-s}}\notag\\
 &\times\exp\left(\sum_{k\ge2}
         \frac{(-1)^{k-1}z^k}{k}\zeta(ks,N+a)\right).
 \label{eq:evaluation}
\end{align}
The tail is convergent when $|z|<(N+a)^{\Rea s}$. The supplied code truncates
this convergent tail and checks consistency with independent formulas;
it does not claim to enclose every floating-point rounding error.
At $s=1$, the removable linear term is evaluated directly as $-\psi(a)$.

\subsection{A confirmed editorial issue, already noticed in incoming work}
In the inspected integration chapter, the historical discussion following
the low-order log-Gamma integrals says: ``Basis atoms must be
$\mathbb Q$-independent.'' Taken literally, this is too strong. One often
cannot prove independence of the retained numerical periods. A valid search
protocol can still remove \emph{known} dependencies and require a proposed
relation to involve its target, then verify or prove it independently.
A suitable replacement is:
\begin{quote}
Remove proved dependencies from the comparison basket, or work in an
explicitly specified quotient. Independence of the retained evaluated
constants need not be known. A returned relation must be checked for its
target coefficient and then subjected to an independent proof or a clearly
labeled numerical validation.
\end{quote}
The incoming shifted-Hurwitz and Stieltjes-calculus reports already propose
a correction of this kind. We confirm the source wording at the pinned
revision and retain the correction as an integration note, not as a newly
discovered flaw. The current pointwise harmonic-number Stieltjes derivative
ladder is not contradicted by anything in this article.

\subsection{Extension guards established by this paper}
Several tempting extensions would be false or underspecified. They are
listed here as guards for integrating \emph{our} formulas, not allegations
that the upstream manuscript made these assertions.

First, $g_d(a)$ is not the diagonal Laurent finite part: already
\eqref{eq:FP2} exhibits the extra $-\gamma_1(a)$. Second, parameter
multiplication needs the exponential in \eqref{eq:multG}; the Gamma slice
is an immediate check. Third, differentiating a moving depth weight is not
fixed-weight spectral differentiation; \eqref{eq:Tclosed} records the
missing logarithmic term. Fourth, a zero of $\G(1,z)$ cannot be substituted
naively into a meromorphic logarithmic derivative: the finite value
\eqref{eq:Froot} results from cancellation.

Fifth, the all-depth convergence domain cannot simply include $s=0$,
as \eqref{eq:s0boundary} shows. Sixth, the nesting bound in
Corollary~\ref{cor:depth} refers to shifted logarithmic Euler sums, not to a
proof of minimal depth or numerical independence in a smaller ordinary
polylogarithm algebra. Finally, the endpoint identity \eqref{eq:endpoint}
uses the convergent largest-index representation, not a sum of undefined
individual all-ones multiple-zeta values.

\subsection{Repository placement and unresolved claims}
The package contains a standalone source, compiled PDF, verification code,
recorded output, a provenance manifest, and integration notes. The incoming
README treats archives as intake material until they are placed and
written into a thematic report; preparation of this delivery does not
perform that intake. No upstream file has been changed.

A natural placement is a new report under the polylogarithm programme,
with cross-references from the integration and differentiation chapters to
Sections~\ref{sec:renorm}--\ref{sec:primitives}, and from the depth chapter
to Sections~\ref{sec:polygen}--\ref{sec:higher}. The target chapter's
numbering, theorem labels, and notation should be reconciled on intake;
the proposed prefix for imported labels is \texttt{nhj:}.

The source already contains an exact $S_4$ proof. This report neither
reproves it as a new result nor resolves the surviving short $S_6$ or newer
$S_8$ candidates. It does not establish new transcendence or numerical
period-independence statements, or claim that all displayed reductions
are first in the literature. Its mathematical contribution is the proved
families above, with their specified regularizations and comparison classes.

\section{Further research questions}
\label{sec:questions}
The identities suggest the following concrete continuations. These are
questions for further work, not additional assertions of this article.

\paragraph{1. Smaller function classes for higher fixed-weight jets.}
Determine whether $U_2^{(m)}$ admits a finite expression using ordinary
polylogarithms, their spectral derivatives, and elementary functions, or
whether a genuinely shifted two-index function is needed in a specified
functional algebra. Equation~\eqref{eq:U22} is an explicit first test.
An obstruction would need a precise differential or symbolic framework;
failure of a numerical search would not establish it.

\paragraph{2. Constructive conversion of the comparison class.}
Give an algorithm converting the rationally shifted logarithmic Euler
sums in \eqref{eq:comparisonclass} into a canonical iterated-integral or
Lerch-derivative vocabulary. Track separately the nesting depth, spectral
derivative order, and the finite shift set. Such a result would sharpen
Corollary~\ref{cor:depth} without making an arithmetic minimality claim.

\paragraph{3. Non-diagonal spectral directions.}
Replace the common exponent by several independent perturbation directions.
Which symmetrized mixed jets still admit a finite-product model, and what
is the correct analogue of the Bell zero-factor theorem? A useful first
case is a two-colour periodic assignment of exponents.

\paragraph{4. Path dependence of multiple Laurent finite parts.}
For a genuinely multivariable nested Hurwitz zeta function, compare limits
along $s_j=1+c_j\eps$ with the equal-direction formula
\eqref{eq:regFP}. Identify the dependence on the direction vector and the
additional algebraic conditions needed for a shuffle- or stuffle-compatible
renormalization. This question must be compared carefully with the existing
renormalization literature, rather than treated as unexplored in general.

\paragraph{5. Individual higher-genus reflection formulas.}
The cyclotomic norm \eqref{eq:cyclogamma} is explicit, but it does not
recover each individual $\G^{[m]}(1/m,z;a)$. Find connection or reflection
laws that complement this norm, particularly at rational $a$ and for
$m=2,3$. Specify growth and logarithm normalizations before proposing
uniqueness of a factorization.

\paragraph{6. Mixed conductor and cyclotomic reductions.}
Study which combinations of rational parameter multiplication and rotation
in the depth variable reduce the individual fractional-order spectral jets
to a finite Dirichlet-$L$-derivative vocabulary. Equation~\eqref{eq:cyclojets}
settles one trace, not all character components.

\paragraph{7. Higher spectral values at moving Gamma zeros.}
Construct an explicit finite algorithm for $F_r(-a-m;a)$ for every $r$,
using the isolated zero factor and pole-free logarithmic jets of the
remaining product. Determine the smallest natural Stieltjes--zeta
coordinate set needed for those evaluations. Proposition~\ref{prop:zeroflow}
gives exact compatibility constraints for any proposed formulas.

\paragraph{8. Argument continuation and monodromy.}
Continue $U_r^{(m)}(w)$ from the disk across arcs away from $w=1$ and give
an exact monodromy calculus. The first-order formula is elementary, but
the higher shifted sums may require a larger alphabet. The coefficient
recurrence supplies a normalized local solution from which to start.

\paragraph{9. Differentiated endpoint identities.}
Determine the exact finite parts after applying powers of $w\partial_w$
to the convergent endpoint sums in Corollary~\ref{cor:endpoint}. At the
first nonsummable power, identify the subtraction polynomial explicitly.
This is an identity problem about normalization, not merely a tail bound.

\paragraph{10. Primitive evaluations at arithmetic depth weights.}
Use \eqref{eq:Apolylog} and \eqref{eq:jetdist1} to isolate combinations of
ordinary and Euler primitives at rational or algebraic $z/a$. Which such
combinations reduce to Gamma, polygamma, Dirichlet-$L$ jets, and a finite
set of algebraic-argument polylogarithms? Any claimed reduction should
retain the underlying branch domain.

\paragraph{11. Connection with the remaining Gaussian mixed candidates.}
Investigate deformations with colour and unequal first order that specialize
to the repository's $S_6$ or $S_8$ comparison baskets. The equal-order
resonance theorem does not currently provide those reductions. A concrete
next step is to derive the coloured finite-product or recurrence identity
before making a new relation search.

\paragraph{12. Formal and certified interfaces.}
Formalize Newton extraction, the regulator-conversion recurrence, the
cyclotomic finite-cutoff identity, and the Bell zero-factor theorem before
attempting the analytic limits. Separately, add interval certification to
the numerical tail evaluators, including cancellation near reciprocal-Gamma
zeros. Neither project should be conflated with the currently executed
floating-point checks.

\section*{Conclusion}
\addcontentsline{toc}{section}{Conclusion}
The equal-order nested family has more exact structure than its individual
depth components suggest. A finite product carries the harmonic coordinates,
its pole subtraction produces a Gamma-valued holomorphic family, and its
Abel transform produces the all-depth multiple polylogarithm. Spectral
jets therefore have two effective normal forms: Stieltjes--zeta Bell
coefficients at the singular spectral slice, and logarithmic harmonic Bell
coefficients at a resonant depth weight. The cyclotomic identity supplies a
further exact bridge from fractional spectral orders to ordinary Gamma.

The main cautions are mathematical rather than merely editorial. A diagonal
Laurent finite part is not a harmonic cutoff constant, a moving depth weight
is not a fixed parameter, and bounded nesting in a stated shifted-sum class
is not numerical period independence. Keeping those distinctions explicit
allows the new formulas to be integrated without reviving the normalization
and evidence problems that the source manuscript itself works to avoid.

\begin{thebibliography}{99}
\bibitem{ProveIt}
V.~Reshetnikov and ProveIt contributors,
\emph{Polylogarithms programme}, repository manuscript and
incoming reports, snapshot \pin, inspected 10 October 2026.
\url{https://github.com/VladimirReshetnikov/ProveIt/tree/fc4d3bf80534ad7c901d3b8c9e71baf2df0064ed/Analysis/Polylogarithms/docs/manuscript}.
Incoming inventory:
\url{https://github.com/VladimirReshetnikov/ProveIt/tree/fc4d3bf80534ad7c901d3b8c9e71baf2df0064ed/docs/incoming}.

\bibitem{DLMFgamma}
NIST Digital Library of Mathematical Functions,
\S5.8, \emph{Infinite Products}, accessed 10 October 2026.
\url{https://dlmf.nist.gov/5.8}.

\bibitem{DLMFmult}
NIST Digital Library of Mathematical Functions,
\S5.5, \emph{Functional Relations}, accessed 10 October 2026.
\url{https://dlmf.nist.gov/5.5}.

\bibitem{DLMFzeta}
NIST Digital Library of Mathematical Functions,
\S25.11, \emph{Hurwitz Zeta Function}, especially its definition,
Euler-summation representations, parameter recurrence, and derivatives,
accessed 10 October 2026.
\url{https://dlmf.nist.gov/25.11}.

\bibitem{DLMFpolylog}
NIST Digital Library of Mathematical Functions,
\S25.12, \emph{Polylogarithms}, accessed 10 October 2026.
\url{https://dlmf.nist.gov/25.12}.

\bibitem{Hoffman2008}
M.~E.~Hoffman,
\emph{A Character on the Quasi-Symmetric Functions coming from Multiple Zeta Values},
Electronic Journal of Combinatorics \textbf{15}(1) (2008), R97.
\url{https://doi.org/10.37236/821}.

\bibitem{Hoffman2016}
M.~E.~Hoffman,
\emph{Harmonic-Number Summation Identities, Symmetric Functions, and Multiple Zeta Values},
arXiv:1602.03198 (2016).
\url{https://arxiv.org/abs/1602.03198}.

\bibitem{IKZ2006}
K.~Ihara, M.~Kaneko, and D.~Zagier,
\emph{Derivation and double shuffle relations for multiple zeta values},
Compositio Mathematica \textbf{142}(2) (2006), 307--338.
\url{https://doi.org/10.1112/S0010437X0500182X}.

\bibitem{GuoZhang2010}
L.~Guo and B.~Zhang,
\emph{Polylogarithms and multiple zeta values from free Rota--Baxter algebras},
Science China Mathematics \textbf{53} (2010), 2239--2258.
\url{https://doi.org/10.1007/s11425-010-4044-1}.
\end{thebibliography}
\end{document}
